\documentclass[11pt]{amsart}
\usepackage[T1]{fontenc}
\usepackage{lmodern,microtype,amsmath,amssymb,amsthm,mathtools,enumitem}
\usepackage[margin=1.1in]{geometry}
\usepackage[hidelinks]{hyperref}
\usepackage{xurl}

\newtheorem{theorem}{Theorem}[section]
\newtheorem{lemma}[theorem]{Lemma}
\newtheorem{corollary}[theorem]{Corollary}
\newtheorem{proposition}[theorem]{Proposition}
\newtheorem*{conjecture*}{Conjecture}
\theoremstyle{remark}
\newtheorem{remark}[theorem]{Remark}
\numberwithin{equation}{section}

\newcommand{\Pp}{\mathbb P}
\newcommand{\E}{\mathbb E}
\newcommand{\Z}{\mathbb Z}
\newcommand{\cB}{\mathcal B}

\hypersetup{pdftitle={The largest box in a sparse random allocation},pdfauthor={Yitzchak Shmalo}}

\title[The largest box in a sparse random allocation]{The largest box in a sparse random allocation}
\author{Yitzchak Shmalo}
\address{Einstein Institute of Mathematics, The Hebrew University of Jerusalem, Jerusalem, Israel}
\email{yitzchak.shmalo@gmail.com}
\date{12 October 2026}
\subjclass[2020]{60C05, 60F05}
\keywords{Random allocation, balls in boxes, largest box, maximum load, numerical semigroup}

\begin{document}
\begin{abstract}
Let $(Y_1,\dots,Y_n)$ be the random allocation of $m$ balls into $n$ boxes in which an allocation $(y_1,\dots,y_n)$ has probability proportional to $\prod_iw_{y_i}$, for a weight sequence $(w_k)_{k\ge0}$ with $w_0>0$, and let $Y_{(1)}$ be the largest of the $Y_i$. Janson conjectured that $Y_{(1)}=o_p(\log n)$ when $m/n\to0$, provided the generating function of the weights has a positive radius of convergence. We prove an explicit tail bound for $Y_{(1)}$, valid for all $n\ge4$ and $1\le m\le cn$ with $c>0$ depending only on the weights, and deduce that $Y_{(1)}\le C+(1+o(1))\,a\log n/\log(n/m)$ with probability tending to one, where $a$ is the smallest $k\ge1$ with $w_k>0$ and $C$ depends only on the set of $k$ with $w_k>0$. This proves the conjecture. For the weights $w_k=1$ the bound is sharp: $Y_{(1)}\log(n/m)/\log n\to1$ in probability when $\log(n/m)=o(\log n)$.
\end{abstract}
\maketitle

\section{Introduction}\label{sec:intro}

Let $w=(w_k)_{k\ge0}$ be a sequence of nonnegative weights with $w_0>0$ and $w_k>0$ for some $k\ge1$. For integers $m\ge0$ and $n\ge1$ let $\cB_{m,n}$ be the set of $y=(y_1,\dots,y_n)\in\Z_{\ge0}^n$ with $y_1+\dots+y_n=m$, let $w(y)=\prod_{i=1}^nw_{y_i}$, and let $Z(m,n)=\sum_{y\in\cB_{m,n}}w(y)$. When $Z(m,n)>0$, the random allocation $B_{m,n}=(Y_1,\dots,Y_n)$ is the random element of $\cB_{m,n}$ with
\begin{equation}\label{eq:law}
\Pp(B_{m,n}=y)=\frac{w(y)}{Z(m,n)},\qquad y\in\cB_{m,n}.
\end{equation}
We think of $Y_i$ as the number of balls in box $i$; only pairs $(m,n)$ with $Z(m,n)>0$ are considered. The weights $w_k=1/k!$ give the multinomial allocation, in which $m$ balls are thrown independently and uniformly into $n$ boxes, and the weights $w_k=1$ give the uniform distribution on $\cB_{m,n}$. Random allocations of this kind and their relation to simply generated random trees are treated in Janson's survey \cite{J12}, whose notation we follow; see also \cite{KSC78,Pav00}.

Let $\Phi(t)=\sum_kw_kt^k$, let $\rho\in[0,\infty]$ be its radius of convergence, let $\Psi(t)=t\Phi'(t)/\Phi(t)$, and let $\nu=\lim_{t\uparrow\rho}\Psi(t)\le\infty$. Arrange $Y_1,\dots,Y_n$ in decreasing order as $Y_{(1)}\ge Y_{(2)}\ge\dots\ge Y_{(n)}$. Suppose that $n\to\infty$ and $m/n\to\lambda$ with $0<\lambda<\nu$, and let $\tau$ be the solution of $\Psi(\tau)=\lambda$. Janson \cite[Theorem~19.3]{J12}, in a result that he describes as essentially due to Meir and Moon \cite{MM90}, showed that then $\tau<\rho$ and
\[
Y_{(1)}\le\frac{\log n}{\log(\rho/\tau)}+o_p(\log n),
\]
that $Y_{(1)}=o_p(\log n)$ if $\rho=\infty$, and that $Y_{(1)}/\log n\to1/\log(\rho/\tau)$ in probability if moreover $w_k^{1/k}\to1/\rho$. As $\lambda\to0$ we have $\tau\to0$, and the constant $1/\log(\rho/\tau)$ tends to $0$. Janson conjectured that Theorem~19.3 holds also for $\lambda=0$, remarking that this seems almost obvious given the case $\lambda>0$, but that there is no general monotonicity \cite[Section~19.1]{J12}. More generally, he leaves the cases $m/n\to0$ and $m/n\to\infty$ as open problems, and refers to \cite{KSC78,Kol86,Pav00,Kaz02} for examples of such results.\footnote{In the arXiv version of \cite{J12} the conjecture is Conjecture~18.4, and Theorem~19.3 is Theorem~18.3.}

\begin{conjecture*}[{Janson \cite[Conjecture~19.4]{J12}}]
If $\rho>0$ and $m/n\to0$, then $Y_{(1)}=o_p(\log n)$.
\end{conjecture*}

We prove the conjecture with an explicit bound, which is of the order $\log n/\log(n/m)$. The constant depends on the arithmetic of the set of box sizes that have positive weight. Let $D=\{k\ge1:w_k>0\}$, which is nonempty, let $a=\min D$, and let $M$ be the additive monoid generated by $D$, that is, the set of finite sums of elements of $D$, the empty sum $0$ included. For each residue class $R$ modulo $a$ that meets $M$ let $c_R=\min(M\cap R)$, and let $C=\max_Rc_R$, the maximum over the at most $a$ classes that meet $M$. Then $C=0$ if and only if $D\subseteq a\Z$; in particular $C=0$ if $w_1>0$. Finally let
\[
\gamma=\min\Bigl(\{1\}\cup\Bigl\{\frac{w_f}{w_0}:f\in D,\ f\le\max(a,C)\Bigr\}\Bigr)\in(0,1].
\]
If $w_1>0$, then $a=1$, $C=0$ and $\gamma=\min(1,w_1/w_0)$.

\begin{theorem}\label{thm:main}
Let $n\ge4$ and $1\le m\le\gamma n/4$ with $Z(m,n)>0$, and let $x=2m/(\gamma n)$. Then for every $j\ge1$
\begin{equation}\label{eq:point}
\Pp(Y_1=j)\le\frac{w_j}{w_0}\,x^{(j-C)/a}.
\end{equation}
Consequently, if $0<r<\rho$ and $q=x^{1/a}/r<1$, then for every integer $k\ge1$
\begin{equation}\label{eq:tail}
\Pp(Y_{(1)}\ge k)\le\frac{\Phi(r)}{w_0}\,n\Bigl(\frac{\gamma n}{2m}\Bigr)^{C/a}\frac{q^k}{1-q}.
\end{equation}
\end{theorem}

\begin{corollary}\label{cor:conj}
Suppose that $\rho>0$. Let $n\to\infty$ and let $m=m(n)\ge1$ satisfy $m/n\to0$ and $Z(m,n)>0$, and put $L=\log(n/m)$. Then for every $\delta>0$
\[
\Pp\Bigl(Y_{(1)}>C+(1+\delta)\frac{a\log n}{L}\Bigr)\to0.
\]
In particular $Y_{(1)}=O_p\bigl(\log n/\log(n/m)\bigr)=o_p(\log n)$.
\end{corollary}

Since $Y_{(1)}=0$ when $m=0$, Corollary~\ref{cor:conj} proves Janson's conjecture. When $w_1>0$ it says that $Y_{(1)}\le(1+\delta)\log n/\log(n/m)$ with probability tending to one, for every $\delta>0$.

When $\log(n/m)=o(\log n)$, the bound is attained for uniform weights, and the factor $a$ cannot be lowered.

\begin{proposition}\label{prop:sharp}
Let $d\ge1$ be an integer, and let $w_k=1$ if $d$ divides $k$ and $w_k=0$ otherwise. Let $n\to\infty$ and let $m=m(n)\ge1$ be divisible by $d$, with $m/n\to0$ and $\log(n/m)=o(\log n)$. Then
\[
\frac{Y_{(1)}\log(n/m)}{\log n}\to d\quad\text{in probability.}
\]
\end{proposition}

For these weights $a=d$ and $C=0$, so the bound of Corollary~\ref{cor:conj} is $(1+\delta)\,d\log n/\log(n/m)$. The additive constant $C$ cannot be dropped either (Remark~\ref{rem:C}). For other weights the largest box can be smaller than the bound by a constant factor (Remark~\ref{rem:multinomial}), and the leading constant $a/\log(n/m)$ agrees to first order with the constant of \cite[Theorem~19.3]{J12} as $\lambda\to0$ (Remark~\ref{rem:janson}).

The proof of Theorem~\ref{thm:main} compares $Z(m-j,n-1)$ with $Z(m,n)$ directly. From an allocation of $m-j$ balls into $n-1$ boxes one obtains allocations of $m$ balls into $n$ boxes by adding an empty box and placing the $j$ balls into $p$ boxes that were empty, in amounts that are fixed elements of $D$. Counting the allocations obtained in this way, and how often each of them arises, gives $Z(m-j,n-1)/Z(m,n)\le(2m/(\gamma n))^p/w_0$, and the arithmetic of $D$ allows $p\ge(j-C)/a$ (Lemmas~\ref{lem:arith} and~\ref{lem:compare}). The argument uses no limit theorem, and the degeneration $\tau\to0$ does not enter.

\section{The set of box sizes}\label{sec:arith}

\begin{lemma}\label{lem:arith}
Let $j\in M$, and let $c$ be the least element of $M$ that is congruent to $j$ modulo $a$. Then $c\le j$, $c\le C$, $a$ divides $j-c$, and $j$ is the sum of $p$ elements of $D$, each at most $\max(a,C)$, where $p\ge(j-C)/a$. Moreover, an integer $j\ge0$ lies in $M$ if and only if its residue class modulo $a$ meets $M$ and $j$ is at least the least element of $M$ in that class.
\end{lemma}

\begin{proof}
Let $R$ be the residue class of $j$ modulo $a$. Then $c=c_R\le j$ and $c\le C$ by the definitions, and $a$ divides $j-c$. Write $c=d_1+\dots+d_t$ with $d_i\in D$, where $t=0$ if $c=0$; each $d_i$ is at most $c\le C$. Let $s=(j-c)/a$. Then $j=d_1+\dots+d_t+sa$ is the sum of $p=t+s$ elements of $D$, each at most $\max(a,C)$, and $p\ge s=(j-c)/a\ge(j-C)/a$. For the last statement, an element of $M\cap R$ is at least $c_R$ by definition, and conversely, if $R$ meets $M$ and $j\in R$ satisfies $j\ge c_R$, then $j=c_R+sa$ with $s\ge0$, and $j\in M$ because $a\in M$ and $M$ is closed under addition.
\end{proof}

\section{A comparison of partition functions}\label{sec:compare}

\begin{lemma}\label{lem:exch}
Let $n\ge2$, let $Z(m,n)>0$ and let $0\le j\le m$. Then for every $1\le i\le n$
\begin{equation}\label{eq:exch}
\Pp(Y_i=j)=\frac{w_j\,Z(m-j,n-1)}{Z(m,n)}.
\end{equation}
\end{lemma}

\begin{proof}
Deleting the $i$-th coordinate is a bijection from $\{y\in\cB_{m,n}:y_i=j\}$ onto $\cB_{m-j,n-1}$, and the weight of each $y$ in this set is $w_j$ times the weight of its image. Summing \eqref{eq:law} over the set gives \eqref{eq:exch}; see also \cite[(18.11)]{J12}.
\end{proof}

In particular the law of $Y_i$ does not depend on $i$.

\begin{lemma}\label{lem:compare}
Let $n\ge2$ and $1\le j\le m$, and let $g_1\le\dots\le g_p$ be elements of $D$ with $g_1+\dots+g_p=j$, each at most $\max(a,C)$. If $p\le n-1-m$, then
\begin{equation}\label{eq:compare}
\binom{n-1-m}{p}\,w_0\gamma^p\,Z(m-j,n-1)\le\binom mp\,Z(m,n).
\end{equation}
\end{lemma}

\begin{proof}
For $y\in\cB_{m-j,n-1}$ let $E(y)=\{i\le n-1:y_i=0\}$. Since $y$ has at most $m-j$ nonempty boxes, $|E(y)|\ge n-1-m\ge p$. For a $p$-element set $J=\{i_1<\dots<i_p\}\subseteq E(y)$ define $z=\phi(y,J)\in\cB_{m,n}$ by $z_{i_s}=g_s$ for $1\le s\le p$, $z_i=y_i$ for $i\le n-1$ with $i\notin J$, and $z_n=0$. The boxes $i_s$ have weight $w_0$ in $y$ and $w_{g_s}$ in $z$, and the added box $n$ has weight $w_0$, so by the definition of $\gamma$
\[
w(\phi(y,J))=w(y)\,w_0\prod_{s=1}^p\frac{w_{g_s}}{w_0}\ge w(y)\,w_0\gamma^p.
\]
Each $y$ belongs to $\binom{|E(y)|}{p}\ge\binom{n-1-m}{p}$ pairs $(y,J)$. On the other hand, if $z=\phi(y,J)$, then $J$ is a $p$-element subset of $\{i\le n-1:z_i\ge1\}$, a set with at most $m$ elements, and $y$ is determined by $z$ and $J$, since $y_i=0$ for $i\in J$ and $y_i=z_i$ for the other $i\le n-1$. Hence each $z\in\cB_{m,n}$ equals $\phi(y,J)$ for at most $\binom mp$ pairs $(y,J)$. Summing $w(\phi(y,J))$ over all pairs therefore gives
\[
\binom{n-1-m}{p}w_0\gamma^p\sum_{y\in\cB_{m-j,n-1}}w(y)\le\sum_{(y,J)}w(\phi(y,J))\le\binom mp\sum_{z\in\cB_{m,n}}w(z),
\]
which is \eqref{eq:compare}.
\end{proof}

\begin{proof}[Proof of Theorem~\ref{thm:main}]
Let $N=n-1-m$. Since $\gamma\le1$ and $n\ge4$, we have $m\le n/4$ and $N\ge3n/4-1\ge n/2\ge m$. Let $j\ge1$. If $j>m$, then $\Pp(Y_1=j)=0$; if $j\le m$ and $w_j=0$, then $\Pp(Y_1=j)=0$ by Lemma~\ref{lem:exch}. In both cases \eqref{eq:point} holds. Otherwise $j\in D\subseteq M$, and Lemma~\ref{lem:arith} writes $j$ as the sum of $p\ge(j-C)/a$ elements of $D$, each at most $\max(a,C)$. As $p\le j\le m\le N$, Lemma~\ref{lem:compare} applies, and together with Lemma~\ref{lem:exch} it gives
\[
\Pp(Y_1=j)=\frac{w_jZ(m-j,n-1)}{Z(m,n)}\le\frac{w_j}{w_0\gamma^p}\,\frac{\binom mp}{\binom Np}=\frac{w_j}{w_0\gamma^p}\prod_{i=0}^{p-1}\frac{m-i}{N-i}\le\frac{w_j}{w_0}\Bigl(\frac{m}{\gamma N}\Bigr)^p,
\]
where we used $(m-i)/(N-i)\le m/N$, which holds because $m\le N$. Since $N\ge n/2$,
\begin{equation}\label{eq:pointp}
\Pp(Y_1=j)\le\frac{w_j}{w_0}\,x^p.
\end{equation}
Finally $x\le1/2<1$ and $p\ge(j-C)/a$, so $x^p\le x^{(j-C)/a}$, and \eqref{eq:point} follows.

For \eqref{eq:tail}, note that $w_j\le\Phi(r)r^{-j}$ for every $j$, so that \eqref{eq:point} gives
\[
\Pp(Y_1=j)\le\frac{\Phi(r)}{w_0}\,r^{-j}x^{(j-C)/a}=\frac{\Phi(r)}{w_0}\,x^{-C/a}q^j .
\]
Since the $Y_i$ have the same law, the union bound gives
\[
\Pp(Y_{(1)}\ge k)\le n\,\Pp(Y_1\ge k)\le\frac{\Phi(r)}{w_0}\,n\,x^{-C/a}\sum_{j\ge k}q^j=\frac{\Phi(r)}{w_0}\,n\,x^{-C/a}\frac{q^k}{1-q},
\]
which is \eqref{eq:tail}, because $x^{-C/a}=(\gamma n/(2m))^{C/a}$.
\end{proof}

\section{Proof of Corollary~\ref{cor:conj}}\label{sec:cor}

Fix $\delta>0$ and $r$ with $0<r<\rho$, and let $b=|\log(\gamma/2)|+a|\log r|$. As $m\ge1$ and $m/n\to0$, we have $L\to\infty$, and for all large $n$ also $n\ge4$ and $m\le\gamma n/4$, so that Theorem~\ref{thm:main} applies with $x=2m/(\gamma n)$. Since $\log(1/x)=L+\log(\gamma/2)$ and $\log(1/q)=(\log(1/x)+a\log r)/a$,
\[
\log(1/x)\le L,\qquad\log(1/q)\ge(L-b)/a .
\]
For all large $n$ we have $L\ge b+a\log2$, so that $q\le1/2$. Let $t=C+(1+\delta)a\log n/L$ and $k=\lfloor t\rfloor+1$, so that $\{Y_{(1)}>t\}=\{Y_{(1)}\ge k\}$ and $k>t>0$. By \eqref{eq:tail} and $1-q\ge1/2$,
\begin{align*}
\log\Pp(Y_{(1)}>t)&\le\log\frac{2\Phi(r)}{w_0}+\log n+\frac Ca\log\frac1x-k\log\frac1q\\
&\le\log\frac{2\Phi(r)}{w_0}+\log n+\frac CaL-\frac{t(L-b)}a\\
&=\log\frac{2\Phi(r)}{w_0}+\frac{Cb}a-\Bigl((1+\delta)\Bigl(1-\frac bL\Bigr)-1\Bigr)\log n .
\end{align*}
The factor in front of $\log n$ tends to $\delta>0$, and $\log n\to\infty$, so $\Pp(Y_{(1)}>t)\to0$. Since $m\ge1$ we have $\log n/L\ge1$, so with probability tending to one $Y_{(1)}\le(C+(1+\delta)a)\log n/L$. Thus $Y_{(1)}=O_p(\log n/L)$, and $Y_{(1)}=o_p(\log n)$ because $L\to\infty$.\qed

\section{Uniform weights}\label{sec:uniform}

\begin{lemma}\label{lem:be}
Let $w_k=1$ for all $k$, and let $n\ge2$ and $1\le k\le m$. Then
\begin{enumerate}[label=\textup{(\roman*)}]
\item $\Pp(Y_1\ge k)=\prod_{i=0}^{k-1}\frac{m-i}{m+n-1-i}$;
\item $\Pp(Y_1\ge k,\ Y_2\ge k)\le\Pp(Y_1\ge k)^2$;
\item $\Pp(Y_{(1)}<k)\le1/\bigl(n\,\Pp(Y_1\ge k)\bigr)$.
\end{enumerate}
\end{lemma}

\begin{proof}
The set $\cB_{m,n}$ has $\binom{m+n-1}{n-1}$ elements, and subtracting $k$ from $y_1$ is a bijection from $\{y\in\cB_{m,n}:y_1\ge k\}$ onto $\cB_{m-k,n}$. Hence $\Pp(Y_1\ge k)=\binom{m-k+n-1}{n-1}/\binom{m+n-1}{n-1}$, which is the product in (i). In the same way $\Pp(Y_1\ge k,\ Y_2\ge k)=\prod_{i=0}^{2k-1}\frac{m-i}{m+n-1-i}$ if $2k\le m$, and the probability is $0$ if $2k>m$. Since $u\mapsto u/(u+n-1)$ is increasing on $u\ge0$, the factor with index $i+k$ is at most the factor with index $i$ for $0\le i<k$, and (ii) follows. For (iii), let $X$ be the number of $i$ with $Y_i\ge k$ and let $\mu=\E X=n\,\Pp(Y_1\ge k)$, which is positive by (i). By symmetry and (ii), $\E X^2=\mu+n(n-1)\Pp(Y_1\ge k,\ Y_2\ge k)\le\mu+\mu^2$, so $\operatorname{Var}X\le\mu$, and by Chebyshev's inequality $\Pp(Y_{(1)}<k)=\Pp(X=0)\le\operatorname{Var}X/\mu^2\le1/\mu$.
\end{proof}

\begin{proof}[Proof of Proposition~\ref{prop:sharp}]
Let $L=\log(n/m)$, so that $L\to\infty$ and $L=o(\log n)$. Consider first $d=1$, and let $0<\delta<1/2$. Here $a=1$, $C=0$ and $\rho=1$, so by Corollary~\ref{cor:conj} $\Pp(Y_{(1)}>(1+\delta)\log n/L)\to0$. For the lower bound let $k=\lfloor(1-\delta)\log n/L\rfloor$. Since $L=o(\log n)$ we have $k\to\infty$ and $m=ne^{-L}=n^{1-o(1)}$; so for all large $n$ we have $L\ge1$, $1\le k\le\log n\le m/2$, and $k/m\to0$. Every factor in Lemma~\ref{lem:be}(i) is at least $(m-k)/(m+n)$, so with $\mu=n\,\Pp(Y_1\ge k)$,
\[
\log\mu\ge\log n-k\log\frac{m+n}{m-k}=\log n-k\Bigl(L+\log\Bigl(1+\frac mn\Bigr)-\log\Bigl(1-\frac km\Bigr)\Bigr)\ge\log n-kL-k\Bigl(\frac mn+\frac{2k}m\Bigr),
\]
using $\log(1+u)\le u$ and $-\log(1-u)\le2u$ for $0\le u\le1/2$. Here $kL\le(1-\delta)\log n$, and $k(m/n+2k/m)\le\log n\,(m/n+2k/m)=o(\log n)$. Hence $\log\mu\ge(\delta-o(1))\log n\to\infty$, and Lemma~\ref{lem:be}(iii) gives $\Pp(Y_{(1)}<k)\le1/\mu\to0$. As $kL/\log n\ge1-\delta-L/\log n$ and $L/\log n\to0$, it follows that $\Pp(Y_{(1)}L/\log n<1-2\delta)\to0$, which together with the upper bound proves the case $d=1$.

For general $d$, $w(y)=1$ if every $y_i$ is divisible by $d$ and $w(y)=0$ otherwise. Hence $B_{m,n}$ is uniformly distributed on the allocations with all $y_i\in d\Z$, and $B_{m,n}/d$ is the allocation with weights $w_k=1$ of $m/d\ge1$ balls into $n$ boxes. Since $\log(n/(m/d))=L+\log d$, the case $d=1$ applies to $m/d$ and gives $(Y_{(1)}/d)(L+\log d)/\log n\to1$ in probability, and $(L+\log d)/L\to1$.
\end{proof}

\section{Remarks}\label{sec:remarks}

\begin{remark}\label{rem:janson}
Let $\rho<\infty$. As $t\to0$ we have $\Phi(t)=w_0+w_at^a+O(t^{a+1})$ and $t\Phi'(t)=aw_at^a+O(t^{a+1})$, so $\Psi(t)=(aw_a/w_0)t^a(1+O(t))$. Hence the solution $\tau$ of $\Psi(\tau)=\lambda$ satisfies $\log(1/\tau)=a^{-1}\log(1/\lambda)+O(1)$ as $\lambda\to0$, and the coefficient $1/\log(\rho/\tau)$ of $\log n$ in \cite[Theorem~19.3]{J12} is $a/\log(1/\lambda)+O\bigl(1/\log^2(1/\lambda)\bigr)$. With $\lambda=m/n$ this is the coefficient $a/\log(n/m)$ of Corollary~\ref{cor:conj}.
\end{remark}

\begin{remark}\label{rem:C}
Let $w_1=0$ and $w_k=1$ for $k\ne1$. Then $D=\{2,3,\dots\}$, $a=2$, $M=\{0,2,3,4,\dots\}$ and $C=3$. If $m=3$, the only allocations of positive weight put all three balls into one box, so $Y_{(1)}=3$, whereas $a\log n/\log(n/m)\to2$. Thus the additive constant $C$ in Corollary~\ref{cor:conj} cannot be omitted.
\end{remark}

\begin{remark}\label{rem:multinomial}
For the multinomial allocation, $w_k=1/k!$, the bound of Corollary~\ref{cor:conj} can be too large by a constant factor. Here $\Pp(Y_1\ge k)\le\binom mkn^{-k}\le(m/n)^k/k!$. Let $\beta>0$ be fixed, let $\log(n/m)=\beta\log\log n+O(1)$ (for instance $m=\lceil n(\log n)^{-\beta}\rceil$), let $\varepsilon>0$ and let $k=\lceil(1+\varepsilon)\log n/((1+\beta)\log\log n)\rceil$. Then $\log k=(1-o(1))\log\log n$ and $k=o(\log n)$, and since $k!\ge(k/e)^k$,
\[
\log\bigl(n(m/n)^k/k!\bigr)\le\log n-\beta k\log\log n-k\log k+O(k)\le-(\varepsilon-o(1))\log n\to-\infty .
\]
By the union bound, with probability tending to one $Y_{(1)}<k$, and hence $Y_{(1)}\le(1+\varepsilon)\log n/((1+\beta)\log\log n)$, while $\log n/\log(n/m)=(1+o(1))\log n/(\beta\log\log n)$. Precise results on the largest box of the multinomial allocation are given in \cite{KSC78,RS98}.
\end{remark}

\section{Computer checks}\label{sec:checks}

The script \texttt{code/checks.py} in the repository listed under Data availability checks the statements of Sections~\ref{sec:arith}--\ref{sec:uniform} in exact rational arithmetic; it uses only the Python standard library, and its output is \texttt{code/checks.out}. It takes eight weight sequences: $w_k=1/k!$; $w_k=1$; $w_k=2^{-k}$; $w_k=k+1$; $w_0=3$ and $w_k=1$ for $k\ge1$; $w_1=0$ and $w_k=1$ otherwise ($a=2$, $C=3$); $w_0=w_3=w_5=1$ and $w_k=0$ otherwise ($a=3$, $C=10$); and $w_0=w_2=1$, $w_4=\frac12$ and $w_k=0$ otherwise ($a=2$, $C=0$). For each of them it checks the membership criterion and the representation of Lemma~\ref{lem:arith} for every $j\le76$; it enumerates all allocations with $n\le7$ and $m\le6$ and checks the weight identity, the numbers of images and preimages and the inequality \eqref{eq:compare} of Lemma~\ref{lem:compare} for every $j\le m$ in $M$ with $p\le n-1-m$; and it computes $Z(m,n)$ and the law of $Y_{(1)}$ exactly by power series, for every $4\le n\le64$ and every $1\le m\le\gamma n/4$ with $Z(m,n)>0$, and for $n\in\{100,200,400\}$ and those $m\in\{1,2,3,5,10,20\}$ with $m\le\gamma n/4$. At each of these points it checks \eqref{eq:pointp} and \eqref{eq:point} for every $j\le m$, and \eqref{eq:tail} for every $k$ with $\max(1,C)\le k\le m$, at one fixed $r$ for each weight sequence, in the stronger form in which $\Phi(r)$ is replaced by $\max_{1\le j\le m}w_jr^j$ and $x^{1/a}$ by a rational lower bound. All $39{,}151$ inequalities hold, and the largest ratio of $\Pp(Y_1=j)$ to the right-hand side of \eqref{eq:point} is $0.5000$ to four decimal places. For $w_k=1$ the script checks Lemma~\ref{lem:be}(i) and the corresponding product formula for $\Pp(Y_1\ge k,\ Y_2\ge k)$ against enumeration for $n\le6$ and $m\le10$, and (ii) and (iii) exactly for $n\le120$, respectively $n\le60$, and $m\le40$. As controls, the script confirms that \eqref{eq:pointp} fails when $p$ is replaced by $p+1$ ($1{,}837$ of $13{,}575$ cases), that \eqref{eq:tail} fails without the factor $n$ ($1{,}032$ of $12{,}001$ cases), and that \eqref{eq:point} fails for the weights of Remark~\ref{rem:C} with $m=3$ when $(j-C)/a$ is replaced by $j/a$. These computations test the identities and inequalities used in the proofs; they are not part of them.

\section*{Acknowledgements and funding}
The research presented in this paper was supported by the European Research
Council (ERC) under the European Union's Horizon Europe research and
innovation programme (grant agreement No.~101041711), by the Simons
Foundation as part of the Collaboration on the Mathematical and Scientific
Foundations of Deep Learning, by Heights Labs, by Convex Nexus Capital, by
the Israel Science Foundation (grant number 2258/19), and by the Israel
Science Foundation (ISF Grant 4101/25). Computational resources were
provided in part by Google Cloud through its education credits program
(credit 530550103).

\section*{Declaration of competing interest}
The author holds equity in Heights Labs and in Convex Nexus Capital and
has a management role at Convex Nexus Capital. These entities are included
in the funding declaration above.

\section*{Declaration of generative AI and AI-assisted technologies}
The paper comes out of Math Brain, an ongoing project on the use of AI systems for research mathematics. Conjecture~19.4 of \cite{J12} was one of the open problems in probability on which the project's agents worked in September 2026. On 9 September 2026 Claude Fable~5.1 (Anthropic) found the proof of Theorem~\ref{thm:main} and Corollary~\ref{cor:conj} in the case $w_1>0$. On 12 October 2026 Claude Opus~5.5 (Anthropic), working in Claude Code, extended the argument to all weight sequences, proved Proposition~\ref{prop:sharp}, searched the literature, and wrote the computer checks and the manuscript. The author directed the work and is responsible for the content of the paper and for decisions concerning its dissemination.

\section*{Data availability}
No data were used. The code is available at\\ \url{https://github.com/yspennstate/sparse-allocation-largest-box}.

\begin{thebibliography}{9}

\bibitem{J12}
S.~Janson, Simply generated trees, conditioned Galton--Watson trees, random allocations and condensation, \emph{Probab. Surv.} \textbf{9} (2012), 103--252.

\bibitem{Kaz02}
N.~I. Kazimirov, On some conditions for absence of a giant component in the generalized allocation scheme, \emph{Discrete Math. Appl.} \textbf{12} (2002), no.~3, 291--302.

\bibitem{Kol86}
V.~F. Kolchin, \emph{Random Mappings}, Optimization Software, New York, 1986.

\bibitem{KSC78}
V.~F. Kolchin, B.~A. Sevast'yanov and V.~P. Chistyakov, \emph{Random Allocations}, V.~H. Winston \& Sons, Washington, D.C., 1978.

\bibitem{MM90}
A.~Meir and J.~W. Moon, On the maximum out-degree in random trees, \emph{Australas. J. Combin.} \textbf{2} (1990), 147--156.

\bibitem{Pav00}
Yu.~L. Pavlov, \emph{Random Forests}, VSP, Zeist, 2000.

\bibitem{RS98}
M.~Raab and A.~Steger, ``Balls into bins''---a simple and tight analysis, in: \emph{Randomization and Approximation Techniques in Computer Science (RANDOM'98)}, Lecture Notes in Comput. Sci. \textbf{1518}, Springer, Berlin, 1998, 159--170.

\end{thebibliography}
\end{document}
